\documentclass[track=junior,level=3]{kidscode}
%% outline stars to color in: \stars{n}
\newcommand{\ostar}{\tikz[baseline=-0.6ex]\node[star,star points=5,star point ratio=2.3,draw=black,line width=0.8pt,fill=white,minimum size=0.75cm,inner sep=0pt]{};}
\newcommand{\stars}[1]{\ifcase#1\or\ostar\or\ostar\,\ostar\or\ostar\,\ostar\,\ostar\fi}
%% plan card boxes (the whole plan card fits on one page in both editions)
\makeatletter\ifkc@dyslexia
  \newcommand{\drawheight}{3.4cm}\newcommand{\partheight}{1.15cm}
\else
  \newcommand{\drawheight}{4.6cm}\newcommand{\partheight}{1.35cm}
\fi\makeatother
%% \planbox{height}{label}: a box to draw or write in, with its label in the corner
\newcommand{\planbox}[2]{%
%% alt: An empty box to draw or write in, labeled with its plan-card step.
\begin{tikzpicture}
  \draw[line width=1.2pt,rounded corners=6pt] (0,0) rectangle (\linewidth,-#1);
  \node[anchor=north west,font=\bfseries,fill=white,draw=black,rounded corners=3pt] at (0.25,-0.2) {#2};
\end{tikzpicture}}
\begin{document}
\kctitle{My Big Build}{Project brief, plan card and stars (weeks 9 and 10)}

\begin{grownup}[How to use this brief]
Read it aloud together at the start of week 9. Your child fills in the plan card before building.
Saying, pointing and drawing all count: only a few words are written.
\end{grownup}

\section{Your mission}
Make your own Big Build. Choose Scratch or Python.\par
Plan it in parts. Build one part at a time.\par
In week 10, you show it at your graduation!

\begin{safety}[Nicknames and credit]
Use a nickname, never your full name.
Started from a starter or a friend's idea? Say so on your plan card!
\end{safety}

\section{Choose your path}
\begin{center}\renewcommand{\arraystretch}{1.25}
\begin{tabularx}{\linewidth}{@{}>{\RaggedRight}X>{\RaggedRight\arraybackslash}X@{}}
\toprule
{\Large\faPuzzlePiece}\ \textbf{Path A: Scratch} & {\Large\faPython}\ \textbf{Path B: Python}\\
A game or an art app. & A turtle postcard.\\
\midrule
{\Large\tickbox}\ It starts with the \textbf{green flag}. &
  {\Large\tickbox}\ It starts from \texttt{postcard.py} and \textbf{runs}.\\
{\Large\tickbox}\ A \textbf{repeat inside a repeat}, or a \textbf{My Block}. &
  {\Large\tickbox}\ A \textbf{\texttt{for} loop}.\\
{\Large\tickbox}\ A \textbf{score} or a \textbf{counter}. &
  {\Large\tickbox}\ At least \textbf{two colors}.\\
 & {\Large\tickbox}\ \textbf{One shape} you changed or added.\\
\midrule
\multicolumn{2}{@{}l@{}}{{\Large\tickbox}\ Both paths: a \textbf{plan card}, and a \textbf{partner check}.}\\
\bottomrule
\end{tabularx}
\end{center}

\section{Idea menu}
Pick one, or think of your own!
\begin{center}
%% alt: Six project ideas in a row, each with a picture: a seedling for a flower garden counter, an apple for a catch-the-apple game, a star for star pattern art, a house for a house postcard, a sun for a sunny day postcard and a robot for a robot face postcard. The first three are Scratch ideas (A), the last three Python ideas (B).
\begin{tikzpicture}[idea/.style={draw=black,line width=1pt,rounded corners=5pt,minimum width=2.6cm,minimum height=2.4cm,align=center,text width=2.3cm,font=\footnotesize}]
  \foreach \icon/\word [count=\i] in {\faSeedling/{A: Flower garden counter},\faApple*/{A: Catch the apple},\faStar/{A: Star pattern art}}
    \node[idea,fill=kcorange!12] at (\i*2.85,0) {{\fontsize{22}{26}\selectfont\icon}\\[3pt]\textbf{\word}};
  \foreach \icon/\word [count=\i] in {\faHome/{B: My house postcard},\faSun/{B: Sunny day postcard},\faRobot/{B: Robot face postcard}}
    \node[idea,fill=kcsky!20] at (8.55+\i*2.85,0) {{\fontsize{22}{26}\selectfont\icon}\\[3pt]\textbf{\word}};
\end{tikzpicture}
\end{center}

\clearpage
\section{My plan card (Milestone 1)}
\noindent\textbf{My Big Build is called} \rule{7.5cm}{0.4pt}\par\medskip
\noindent\textbf{Circle one:}\quad {\Large\faPuzzlePiece}\ \textbf{A: Scratch}\qquad {\Large\faPython}\ \textbf{B: Python}\par\medskip

\noindent\planbox{\drawheight}{Draw it: what will it look like?}\par\medskip
\noindent\textbf{My three parts, in order:}\par\smallskip
\noindent\planbox{\partheight}{Part 1}\par\smallskip
\noindent\planbox{\partheight}{Part 2}\par\smallskip
\noindent\planbox{\partheight}{Part 3}\par\medskip

\noindent\textbf{My must-haves} (check them when they are in your plan):\par
\noindent A:\ {\Large\tickbox}\,green flag\quad{\Large\tickbox}\,repeat in a repeat or My Block\quad{\Large\tickbox}\,score or counter\par
\noindent B:\ {\Large\tickbox}\,it runs\quad{\Large\tickbox}\,\texttt{for} loop\quad{\Large\tickbox}\,two colors\quad{\Large\tickbox}\,my shape\par\medskip
\noindent\textbf{Whose idea did I use?} \rule{8cm}{0.4pt}\par\medskip
\noindent\textbf{Partner check by} (nickname) \rule{5cm}{0.4pt}\quad {\Large\tickbox}\ done

\clearpage
\section{Milestones}
Build one part at a time. Test after every change.

\begin{center}\renewcommand{\arraystretch}{1.6}
\begin{tabularx}{\linewidth}{@{}c>{\bfseries}p{4.4cm}X@{}}
\toprule
 & Milestone & Done when \ldots\\
\midrule
\tickbox & M1 Plan card\newline{\normalfont\small week 9, start} & I chose A or B. I drew it and named three parts.\\ \midrule
\tickbox & M2 It works\newline{\normalfont\small week 9, end} & It runs. It has every must-have. A partner checked it. I saved it.\\ \midrule
\tickbox & M3 Graduation showcase\newline{\normalfont\small week 10} & I showed my Big Build. I said what one part of my code does.\\
\bottomrule
\end{tabularx}
\end{center}

\section{Need a different path?}
\begin{tip}[Core path: start from a template]
A grown-up gets your project ready. Path A: a Scratch project with a backdrop, the cat, the Pen
and a counter. Path B: your own copy of \texttt{postcard.py}. It already draws a house and a roof!
Change one \texttt{CHANGE ME} line at a time. Press Run after each change.
\end{tip}

\noindent\begin{minipage}{\linewidth}
\begin{pycode}[postcard.py: change the color, then press Run]
t.color("orange")   # CHANGE ME: the color of the house
for i in range(4):
    t.forward(100)
    t.left(90)
\end{pycode}
\end{minipage}

\begin{stretchchallenge}[Stretch path: my own block]
\begin{minipage}[t]{0.5\linewidth}\RaggedRight
Do you draw the same shape again and again?\par\smallskip
Give it a name with \texttt{def}. It is a Python My Block!\par\smallskip
The lines inside are pushed in. Then use it with its name.
\end{minipage}\hfill
\begin{minipage}[t]{0.44\linewidth}\vspace{0pt}
\begin{pycode}
def box():
    for i in range(4):
        t.forward(50)
        t.left(90)

box()
\end{pycode}
\end{minipage}
\end{stretchchallenge}

\needspace{22\baselineskip}
\section{How my Big Build is checked}
Color the stars you have now. Look again after the showcase!

\begin{rubric}[My stars]
\renewcommand{\arraystretch}{1.3}
\begin{tabularx}{\linewidth}{@{}>{\bfseries\RaggedRight}p{3.6cm}*{3}{>{\RaggedRight\arraybackslash}X}@{}}
 & \stars{1} & \stars{2} & \stars{3}\\
\midrule
It works & It starts. & Most parts work. & It all works, every time.\\ \midrule
It has every must-have & one of them & most of them & all of them\\ \midrule
I can explain one part of my code & I can say what my Big Build does. & I can point to a part and name its blocks or lines. & I can say what each block or line in one part does.\\ \midrule
I planned it on my plan card & I filled in a plan card. & My Big Build has some of my plan. & My Big Build matches my three parts.\\
\end{tabularx}
\end{rubric}

\begin{reflection}[After the showcase]
Draw the part of your Big Build you like best. How do you feel? Circle one.
\begin{center}
\Huge\faSmile[regular]\qquad\faMeh[regular]\qquad\faFrown[regular]
\end{center}
\answerspace{2.6cm}
\end{reflection}
\end{document}
